\documentclass[11pt]{article}
\usepackage[margin=1.1in]{geometry}
\usepackage{amsmath,amssymb,amsthm}
\usepackage{booktabs}
\usepackage[expansion=false]{microtype}
\usepackage{xcolor}
\usepackage[colorlinks=true,linkcolor=blue!50!black,citecolor=blue!50!black,urlcolor=blue!50!black]{hyperref}

\newtheorem{theorem}{Theorem}[section]
\newtheorem{lemma}[theorem]{Lemma}
\newtheorem{proposition}[theorem]{Proposition}
\newtheorem{corollary}[theorem]{Corollary}
\theoremstyle{definition}
\newtheorem{problem}{Problem}
\newtheorem{definition}[theorem]{Definition}
\theoremstyle{remark}
\newtheorem{remark}[theorem]{Remark}
\newtheorem{observation}[theorem]{Observation}

\newcommand{\jac}[2]{\left(\tfrac{#1}{#2}\right)}
\DeclareMathOperator{\ord}{ord}

\title{Local Obstructions and Slowly Growing Search Depth\\
in the Erd\H{o}s--Straus Conjecture:\\
A Classified Census to $10^{15}$}
\author{Vincent\thanks{Computations performed with machine assistance; all numerical
claims are reproducible from the released code and data (\S\ref{sec:data-availability}).
Code, data, and discussion are hosted at
\url{https://github.com/Attitude-Adjuster/erdos-strauss}; questions, corrections, and
reproducibility reports may be filed through the repository's issue tracker.}}
\date{August 2026 --- working draft, version 6}

\begin{document}
\maketitle

\begin{abstract}
For a prime $p$, write $A=(p+r)/4$ with $r\equiv3\pmod4$ and $q=pA$.  The
Erd\H{o}s--Straus equation has a solution at this $A$ exactly when a divisor of $q^2$
lies in the class $-q\pmod r$.  This turns the search into a ladder over the rungs
$r=3,7,11,\ldots$ and separates failed rungs into four levels: obstruction by the
Jacobi symbol, by another real character, by the multiplicative support subgroup, or a
\emph{residual miss} for which the target class is locally available but no divisor
reaches it.

We classify every failed rung for the $932{,}642{,}160{,}749$ primes below $10^{15}$ in
the six square classes modulo $840$ not covered by Mordell's identities. Among
$558$ billion failed rungs, $99.1773\%$ have a Jacobi certificate,
$0.00008\%$ require a finer local certificate, and $0.8227\%$ are residual.  The finest
local level occurs only at the single rung $r=51$, in all $43{,}257$ instances, and we prove
that it must: since $4A=p+r$, the support subgroup always contains $4$---a constraint
invisible to every real character, because $\chi(4)=\chi(2)^2=1$, but decisive for the
support test.  It confines that level to ten rungs below $300$, of which $51$ is the
smallest.  The
residual share falls monotonically across the census, but its absolute count continues
to grow.  The maximum first-hit rung also rises, from the long-standing record $107$
through $131$ to $155$.  The data therefore gives no support for a uniform constant bound on search
depth; instead it suggests an unbounded, exceptionally slowly growing envelope.

The residual misses are concentrated among divisor-poor shifted values $(p+r)/4$.
However, a random-occupancy model for the divisors of $q^2$ fails by orders of
magnitude, showing that multiplicative dependence, not merely divisor density, is the
central obstacle.  These results are numerical and do not prove a new case of the
conjecture.  Their contribution is a reproducible decomposition of what local residue
information explains and what it leaves arithmetically unresolved.
\end{abstract}

\section{Introduction}\label{sec:intro}

The Erd\H{o}s--Straus conjecture \cite{Erdos1950} asserts that for every integer
$n \ge 2$ the equation $4/n = 1/x + 1/y + 1/z$ has a solution in positive integers; it
suffices to treat prime $n = p$. Polynomial congruence identities resolve the equation
outside a thin set of residue classes. Mordell's identities leave six classes
$p \equiv 1^2, 11^2, 13^2, 17^2, 19^2, 23^2 \pmod{840}$, all of them squares modulo
$840$ \cite{Mordell1969}. We denote their union by the \emph{hard class}
$\mathcal{H}$. The common surrogate $p\equiv1\pmod{24}$ is four times larger and mixes
these primes with classes already settled by an identity; Section~\ref{sec:hardclass}
shows that this dilutes the measured residual share. We therefore census
$\mathcal{H}$ directly. The obstruction is structural: Schinzel and Yamamoto showed
that no finite list of polynomial congruence identities covers a quadratic-residue
class \cite{Schinzel2000,Yamamoto1965}.

Fixing $A$ and setting $r = 4A - p$, $q = pA$, solvability is equivalent to the existence of a
divisor of $q^2$ in the class $-q \bmod r$ (Lemma~\ref{lem:divisor}, classical).
Failure can then have two very different causes:

\begin{itemize}
\item \emph{Local residue support.} The divisors of $q^2$ all lie in some proper
subset of $(\mathbb{Z}/r)^\times$ determined by the residues of the primes of $q$, and
the target class lies outside it. No search is required; a short certificate settles it.
\item \emph{A residual miss.} The target class is available to the divisors in every
local sense, but no divisor happens to land there.
\end{itemize}

The first is finite group arithmetic. The second is the unresolved part of the search.
We make this separation explicit through a nested certificate hierarchy, classify it
exhaustively on $\mathcal{H}$ below $10^{15}$, and release independently checkable
certificates. The parametrization, character expansion, and Type I/II taxonomy are
classical; the contribution here is the hierarchy, a structural theorem pinning its finest
local level to an explicit finite list of rungs (\S\ref{sec:whyS}), the exhaustive
classification, and what it reveals about the slowly growing residual tail.

\section{The rung ladder and the divisor lattice}\label{sec:ladder}

Fix a prime $p \equiv 1 \pmod 4$ and an integer $A > p/4$, and set
$r = 4A - p > 0$, $q = pA$. Subtracting $1/A$ from $4/p$ gives $1/B + 1/C = r/q$, and
completing the product turns this into a hyperbola:
\begin{equation}\label{eq:hyperbola}
\frac{4}{p} = \frac{1}{A} + \frac{1}{B} + \frac{1}{C}
\qquad\Longleftrightarrow\qquad
(rB - q)(rC - q) \;=\; q^2 .
\end{equation}

\begin{lemma}[divisor criterion; classical]\label{lem:divisor}
Positive integer solutions $(B, C)$ of \eqref{eq:hyperbola} correspond bijectively to
divisors $u \mid q^2$ with $u \equiv -q$ and $q^2/u \equiv -q \pmod r$, via
$B = (u+q)/r$, $C = (q^2/u + q)/r$. If $\gcd(q, r) = 1$ the second congruence follows
from the first.
\end{lemma}

\begin{proof}
If $(B,C)$ solves \eqref{eq:hyperbola} then $u = rB - q$ divides $q^2$ and is
$\equiv -q \pmod r$, as is its cofactor; conversely such a $u$ yields integers
$B, C \ge 1$. When $\gcd(q,r)=1$, $u \equiv -q$ gives
$q^2/u \equiv q^2(-q)^{-1} \equiv -q \pmod r$.
\end{proof}

As $A$ runs over integers exceeding $p/4$, $r$ runs over $r \equiv -p \equiv 3 \pmod 4$.
We call the resulting search order a \emph{ladder} with \emph{rungs}
$r = 3, 7, 11, \dots$ and $A = (p+r)/4$; rung $r$ \emph{hits} if
Lemma~\ref{lem:divisor} produces a solution, and the \emph{depth} of $p$ is the
smallest hitting $r$.

The required coprimality is automatic throughout the range of interest.

\begin{lemma}\label{lem:coprime}
If $0 < r < 3p$ then $\gcd(q, r) = 1$.
\end{lemma}

\begin{proof}
$r < 3p$ is equivalent to $A < p$, so $p \nmid A$ and, $p$ being prime,
$\gcd(A,p) = 1$. Then $\gcd(A, r) = \gcd(A, 4A - p) = \gcd(A, p) = 1$ and
$\gcd(p, r) = \gcd(p, 4A - p) = \gcd(p, 4A) = 1$ since $p$ is odd and $p \nmid A$.
Hence $\gcd(pA, r) = 1$.
\end{proof}

Thus the second congruence in Lemma~\ref{lem:divisor} is not an additional condition in
the census.

The divisors of $q^2 = p^2A^2$ form a box lattice: writing $A = p^{e}A'$ with
$p \nmid A'$ and $A' = \prod_i \ell_i^{a_i}$, they are $u = p^{\,j}d$ with
$0 \le j \le 2e+2$ and $d \mid A'^2$. Reducing modulo $r$ colors this grid, and
Lemma~\ref{lem:divisor} asks whether the color $-q$ appears. The classical two-type
taxonomy is the row coordinate:

\begin{proposition}[types are rows]\label{prop:types}
Suppose $p \nmid A$ and $\gcd(q,r) = 1$, and let a solution arise from
$u = p^{\,j}d$ with $j \in \{0,1,2\}$ and $d \mid A^2$. If $j = 1$ then $p \mid B$ and
$p \mid C$ \textup{(}Type II\textup{)}; if $j \in \{0,2\}$ then $p$ divides exactly one
of $B, C$ \textup{(}Type I\textup{)}.
\end{proposition}

\begin{proof}
$p \mid q$, so $p \mid u + q$ iff $p \mid u$, i.e.\ iff $j \ge 1$; the cofactor has
$p$-exponent $2-j$. Since $\gcd(r,p)=1$, $p \mid B$ iff $p \mid u+q$.
\end{proof}

\begin{remark}
For $p = 73$, $A = 20$ ($r = 7$, $q = 1460$) the lattice contains four solution
divisors $u \le q$: $u = 10, 80$ in row $j = 0$ (Type I, e.g.\
$4/73 = 1/20 + 1/210 + 1/30660$) and $u = 73, 584$ in row $j = 1$ (Type II, e.g.\
$4/73 = 1/20 + 1/219 + 1/4380$). Both classical families sit in one lattice.
\end{remark}

\section{The character expansion and a hierarchy of certificates}\label{sec:fourier}

Throughout, $\gcd(q,r) = 1$ and
$N_r(p) = \#\{u \mid q^2 : u \equiv -q \ (\mathrm{mod}\ r)\}$, counting all divisors
(each unordered solution $\{B,C\}$ contributes two).

\begin{theorem}[divisor count in a class; orthogonality]\label{thm:fourier}
With $q^2 = \prod_\ell \ell^{\,a_\ell}$,
\begin{equation}\label{eq:fourier}
N_r(p) \;=\; \frac{1}{\varphi(r)} \sum_{\chi \ (\mathrm{mod}\ r)}
\overline{\chi(-q)} \prod_{\ell \mid q}
\bigl( 1 + \chi(\ell) + \chi(\ell)^2 + \cdots + \chi(\ell)^{a_\ell} \bigr).
\end{equation}
\end{theorem}

\begin{proof}
For $u, a$ coprime to $r$, orthogonality \cite[Ch.~4]{Davenport2000} gives
$\varphi(r)^{-1}\sum_\chi \chi(u)\overline{\chi(a)} = \mathbf{1}[u \equiv a]$. Take
$a = -q$, sum over $u \mid q^2$, and factor the completely multiplicative $\chi$ across
the divisor structure of $q^2$.
\end{proof}

The conjugate in \eqref{eq:fourier} is essential: omitting it counts the class
$(-q)^{-1}$ rather than $-q$. Table~\ref{tab:verify} checks both evaluations against
direct enumeration.

The trivial character contributes $\tau(q^2)/\varphi(r) > 0$; the others contribute
cancellation. For real characters that cancellation is itself a divisor count:

\begin{lemma}\label{lem:interference}
Let $\chi$ be a real character mod $r$. Since every exponent $a_\ell$ in $q^2$ is even,
\[
\sum_{u \mid q^2} \chi(u) \;=\; \prod_{\chi(\ell) = 1}(a_\ell + 1) \;=\; \tau\bigl((q^2)_+\bigr) \;\ge\; 0,
\]
where $(q^2)_+$ is the part of $q^2$ supported on primes with $\chi(\ell) = 1$: at
primes with $\chi(\ell) = -1$ the local factor $1 - 1 + \cdots + 1$ equals $1$.
\end{lemma}

The obstruction applies to every odd $r$ and every real character.

\begin{corollary}[real-character certificate]\label{cor:blocking}
Let $r$ be odd with $\gcd(q,r) = 1$ and let $\chi$ be any real Dirichlet character
mod $r$. If $\chi(\ell) = +1$ for every prime $\ell \mid q$ and $\chi(-q) = -1$, then
$N_r(p) = 0$.
\end{corollary}

\begin{proof}
$\chi$ is completely multiplicative, so every divisor $u \mid q^2$ has $\chi(u) = +1$,
whereas $\chi$ is constant on residue classes mod $r$ and equals $-1$ at $-q$. Hence no
divisor lies in that class. Equivalently in \eqref{eq:fourier}: by
Lemma~\ref{lem:interference} the $\chi$ term is $\chi(-q)\tau(q^2) = -\tau(q^2)$,
cancelling the trivial term exactly.
\end{proof}

\begin{remark}[the Jacobi symbol is only one of them]\label{rem:whichchi}
For odd $r = \prod_{i=1}^{\omega} p_i^{e_i}$ the group $(\mathbb{Z}/r)^\times$ has
exactly $2^{\omega}$ real characters, namely the products
$\chi_T(x) = \prod_{i \in T} \jac{x}{p_i}$ over subsets
$T \subseteq \{1,\dots,\omega\}$. The Jacobi symbol is one of them---but \emph{not}
in general the full-subset one: $\jac{\cdot}{r} = \chi_{T_0}$ with
$T_0 = \{i : e_i \text{ odd}\}$, since $\jac{x}{p_i^{e_i}} = \jac{x}{p_i}^{e_i}$. For
squarefree $r$ these coincide; for $r = 63 = 3^2\cdot 7$ the Jacobi symbol is
$\jac{\cdot}{7}$ and for $r = 75 = 3\cdot 5^2$ it is $\jac{\cdot}{3}$, and both of
these rungs occur in our data. For prime $r$ there is a single nontrivial real
character and the distinction is vacuous; for composite $r$, testing only the Jacobi
symbol is strictly weaker than testing all real characters. Note also that
$\jac{x}{r} = +1$ does not mean $x$ is a quadratic residue when $r$ is composite, so
the hypothesis must be phrased in terms of the character, not of residuosity.
\end{remark}

The criteria simplify considerably once one uses $r \equiv 3 \pmod 4$, and the
simplification explains exactly which rungs can exhibit which levels.

\begin{proposition}[criteria in reduced form]\label{prop:reduced}
Let rung $r$ fail for $p$, and let $H \le (\mathbb{Z}/r)^\times$ be as in
Definition~\ref{def:hierarchy}. Recall $r \equiv 3 \pmod 4$.
\begin{enumerate}
\item A real character $\chi$ mod $r$ certifies the failure if and only if
$\chi(-1) = -1$ \textup{(}$\chi$ is odd\textup{)} and $H \le \ker\chi$. In
particular the Jacobi symbol certifies if and only if $\jac{\ell}{r} = +1$ for every
prime $\ell \mid q$: the hypothesis $\jac{-q}{r} = -1$ is automatic.
\item $-q \in H$ if and only if $-1 \in H$. Hence the support certificate of
Definition~\ref{def:hierarchy} says exactly that $-1 \notin H$.
\item There are exactly $2^{\omega(r)-1}$ odd real characters mod $r$, one of which is
the Jacobi symbol. Consequently a level-\textup{RC} failure can occur only when
$\omega(r) \ge 2$.
\end{enumerate}
\end{proposition}

\begin{proof}
Since $q$ is a product of the primes generating $H$, we have $q \in H$; so if
$H \le \ker\chi$ then $\chi(q) = 1$ and $\chi(-q) = \chi(-1)$, giving (1). For the
Jacobi symbol, $\jac{-1}{r} = (-1)^{(r-1)/2} = -1$ because $r \equiv 3 \pmod 4$.
Part (2) is the same observation: $-q \in H \iff -1 \in H$ because $q \in H$. For (3),
$\chi_T(-1) = \prod_{i \in T}(-1)^{(p_i-1)/2}$, so $\chi_T$ is odd exactly when $T$
meets $\{i : p_i \equiv 3 \ (\mathrm{mod}\ 4)\}$ in an odd number of elements; that
set is nonempty since $r \equiv 3 \pmod 4$, and the number of such $T$ is
$2^{\omega-1}$. When $\omega = 1$ there is therefore exactly one odd real character,
which must be the Jacobi symbol, leaving no room for level \textup{RC}.
\end{proof}

\begin{definition}[certificate hierarchy]\label{def:hierarchy}
Let rung $r$ fail for $p$ (so $\gcd(q,r) = 1$ by Lemma~\ref{lem:coprime}), and let
$H \le (\mathbb{Z}/r)^\times$ be the subgroup generated by
$\{\ell \bmod r : \ell \mid q\}$. The failure has level
\begin{description}
\item[J] if the Jacobi symbol $\jac{\cdot}{r}$ satisfies Corollary~\ref{cor:blocking};
\item[RC] if some real character mod $r$ does, but the Jacobi symbol does not;
\item[S] if no real character does, but $-q \bmod r \notin H$;
\item[R] (\emph{residual miss}) otherwise: $-q \in H$, yet the divisors of $q^2$---whose
exponents are bounded---miss the class.
\end{description}
\end{definition}

\begin{proposition}\label{prop:nested}
The levels are nested. Every J-certificate is an RC-certificate, and every
RC-certificate implies $-q \notin H$. Level R is what remains after all local support
information is exhausted.
\end{proposition}

\begin{proof}
The first claim is Remark~\ref{rem:whichchi}. For the second, if $\chi$ is trivial on
every $\ell \mid q$ then $H \le \ker\chi$, while $\chi(-q) = -1$ puts $-q$ outside
$\ker\chi$.
\end{proof}

Levels RC and S are finer tests in principle---S detects obstructions from
higher-order characters and from subgroups of index greater than $2$ that no real
character sees. Whether they are finer \emph{in practice}, and at what height they
first bite, is an empirical question answered in \S\ref{sec:data};
\S\ref{sec:whyS} shows that level S is in addition confined to an explicit finite set of
rungs.

For the smallest rung, level R is empty unconditionally, because modulo $3$ there are
only two characters.

\begin{proposition}[rung 3 has no residual misses]\label{prop:rung3}
Let $p \equiv 1 \pmod{12}$ and $A = (p+3)/4$, so $r = 3$. Then rung $3$ hits unless it
carries a real-character certificate; explicitly, it fails precisely when $3 \nmid A$
and every prime $\ell \mid pA$ satisfies $\ell \equiv 1 \pmod 3$.
\end{proposition}

\begin{proof}
If $3 \mid A$ take $u = 3$: then $3 \mid q$, $9 \mid q^2$, and both congruences of
Lemma~\ref{lem:divisor} read $0 \equiv 0 \pmod 3$, so rung $3$ hits. If $3 \nmid A$
then $\gcd(q,3) = 1$ and \eqref{eq:fourier} has two terms,
$N_3(p) = \tfrac12\bigl(\tau(q^2) + \chi(-q)\,\tau((q^2)_+)\bigr)$ with $\chi$ the
nontrivial character mod $3$ (real, so conjugation is immaterial). If some $\ell \mid q$
has $\chi(\ell) = -1$ then $\tau((q^2)_+) < \tau(q^2)$ and $N_3 > 0$ either way. If
every $\ell \mid q$ has $\chi(\ell) = 1$ then $\tau((q^2)_+) = \tau(q^2)$ and
$\chi(-q) = \chi(-1) = -1$, so $N_3 = 0$; since $\jac{\cdot}{3}$ is the only nontrivial
real character mod $3$, this is a level-J certificate.
\end{proof}

\subsection{Which rungs can carry a level-S failure}\label{sec:whyS}

Level S is the one level no real-character certificate detects, and empirically
(\S\ref{sec:data}) it sits at a single rung. Higher-order characters still detect the
support obstruction, as the centered Fourier formulation in \S\ref{sec:doorways} makes
explicit. Its confinement is forced by a constraint on $H$ that every real character is
blind to.

\begin{lemma}[the four-constraint]\label{lem:four}
$4 \in H$, always.
\end{lemma}

\begin{proof}
$4A = p + r$ gives $A \equiv 4^{-1}p \pmod r$, the inverse existing because $r$ is odd.
Now $H$ contains $p$ and every prime factor of $A$, hence contains $A$ itself; therefore
$4^{-1} = A\,p^{-1} \in H$ and so $4 \in H$.
\end{proof}

Any real $\chi$ has $\chi(4) = \chi(2)^2 = 1$, so $4 \in \ker\chi$ automatically and
Lemma~\ref{lem:four} excludes nothing in Corollary~\ref{cor:blocking}. For the support
test it is decisive.

\begin{proposition}[when level S is possible]\label{prop:Spossible}
Write $G = (\mathbb{Z}/r)^\times$. Rung $r$ can carry a level-S failure only if some
subgroup $H \le G$ satisfies
\begin{equation}\label{eq:Sadm}
4 \in H, \qquad -1 \notin H, \qquad (-1)H \in (G/H)^2 ,
\end{equation}
and any $H$ satisfying \eqref{eq:Sadm} does produce one whenever it is realized as a
support subgroup. In particular level S is impossible at every prime rung.
\end{proposition}

\begin{proof}
The first condition is Lemma~\ref{lem:four}, the second is
Proposition~\ref{prop:reduced}(2), and the second also forces the rung to fail, since all
divisors of $q^2$ lie in $H$. For the third: the real characters trivial on $H$ are
exactly the characters of $G/H$ of order dividing $2$, and by
Proposition~\ref{prop:reduced}(1) the failure escapes every real-character certificate
precisely when all of them take the value $+1$ at $-1$. The intersection of the kernels of
all order-$\le 2$ characters of a finite abelian group is its subgroup of squares, which
is the stated condition.

If $r$ is prime then $G$ is cyclic of order $r - 1 \equiv 2 \pmod 4$. Then $-1 \notin H$
forces $|H|$ odd, so $G/H$ is cyclic of order $2m$ with $m$ odd and $(-1)H$ is its unique
element of order $2$ --- not a square, since $4 \nmid |G/H|$. No $H$ qualifies.
\end{proof}

\begin{corollary}\label{cor:Srungs}
Among $r \equiv 3 \pmod 4$ with $r < 300$, level S is possible only for
\[
r \;\in\; \{51,\ 119,\ 123,\ 187,\ 195,\ 219,\ 255,\ 267,\ 287,\ 291\},
\]
and $51$ is the smallest by a wide margin. At $r = 51$, $119$ and $123$ the admissible
subgroup is unique and cyclic of index $4$; at $r = 51$ it is
$H = \langle 2 \rangle = \{1,2,4,8,13,16,26,32\}$ inside
$(\mathbb{Z}/51)^\times \cong \mathbb{Z}/2 \times \mathbb{Z}/16$.
\end{corollary}

Dropping Lemma~\ref{lem:four} from \eqref{eq:Sadm} would admit fourteen rungs below $140$
alone,
\[
15,\ 35,\ 39,\ 51,\ 55,\ 75,\ 87,\ 91,\ 95,\ 111,\ 115,\ 119,\ 123,\ 135,
\]
with nothing to distinguish $51$, and with $r = 15$ visited far more often than $r = 51$. The
four-constraint eliminates eleven of the fourteen. It also disposes of the objection that
the confinement might be an artifact of which rungs are reached: level RC is observed at
$r = 51, 75, 91, 99, 115$, so those rungs \emph{are} reached carrying local obstructions,
yet at $75$, $91$ and $115$ every subgroup that would otherwise satisfy \eqref{eq:Sadm}
omits $4$, and at $99$ none exists in the first place.

The statement is falsifiable at fixed cost: \textbf{no census, at any height, may report a
level-S failure at $r = 15, 35, 39, 55, 75, 87, 91, 95, 111, 115$ or $135$}, and the next
rung that can host one is $r = 119$.

\section{Exhaustive classification below \texorpdfstring{$10^{15}$}{1e15}}\label{sec:data}

The census follows every prime in $\mathcal{H}$ to its first hitting rung and classifies
all preceding failures by Definition~\ref{def:hierarchy}. It contains
$932{,}642{,}160{,}749$ primes and $557{,}964{,}077{,}963$ failed rungs below $10^{15}$;
the $10^{13}$ stage was produced on CPUs and the extension to $10^{15}$ by a GPU
implementation whose statistics and certificate set are byte-identical to the CPU
reference on every range where both run. Independent exact-arithmetic and machine-word
implementations agree on all overlapping ranges, and every emitted fine-certificate or
deep-prime claim---all $585{,}677$ of them---was checked again in exact arithmetic on a
machine different from the one that produced it. Section~\ref{sec:computation} gives
the algorithm and verification protocol.

\begin{table}[ht]\centering
\caption{Observed maximum search depth $D(p)$, the first hitting rung, on
$\mathcal{H}$. Long plateaus are followed by small record increments; they are evidence
of an extreme tail, not of a fixed ceiling.}\label{tab:depth}
\begin{tabular}{lrr}
\toprule
cutoff $X$ & $\max_{p<X}D(p)$ & first record holder in the census \\
\midrule
$10^8$    & $107$ & $8{,}803{,}369$ \\
$10^{10}$ & $107$ & $8{,}803{,}369$ \\
$10^{11}$ & $107$ & $8{,}803{,}369$ \\
$10^{12}$ & $111$ & $654{,}730{,}707{,}409$ \\
$10^{13}$ & $131$ & $6{,}367{,}024{,}012{,}441$ \\
$10^{14}$ & $131$ & $6{,}367{,}024{,}012{,}441$ \\
$10^{15}$ & $155$ & $172{,}538{,}390{,}619{,}841$ \\
\bottomrule
\end{tabular}
\end{table}

The depth record remained $107$ across three decades of new primes, then moved to
$111$, to $131$, and---within the fifteenth decade---to $155$, at
\[
p = 172{,}538{,}390{,}619{,}841, \qquad
\begin{aligned}
\frac{4}{p} &= \frac{1}{43134597654999} + \frac{1}{48015316516502146586034030} \\[2pt]
            &\quad + \frac{1}{567275091541152628215802033426900470},
\end{aligned}
\]
and the displayed decomposition is exact. Below $10^{15}$ the depths
$111$, $115$, $119$, $123$, $127$, $131$ (eight primes), $135$ (three), and $139$ (one)
are all occupied, while $143$, $147$, and $151$ are not: the record sits beyond a gap,
as $107$ once did. The long plateaus are therefore sampling effects in an exceptionally
thin tail. Numerically, the natural expectation is not a constant bound $D(p)\le R$,
but an unbounded envelope that grows too slowly to estimate reliably from record values
alone. This is suggestive, not a theorem; the current record is no more a bound than its
predecessor was.

All failures at $r=3$ are level J, as Proposition~\ref{prop:rung3} predicts. Residual
misses first appear at $r=11$ and are concentrated at small rungs, while the finer local
levels occur only on the thin subsequence of composite rungs. Their full counts and
localization are reported below on the same population as the main census.

\subsection{Which primes are actually hard}\label{sec:hardclass}

Mordell's identities \cite{Mordell1969} settle $4/n$ for every $n$ except possibly those
$n$ that are \emph{squares} modulo $840$. There are exactly six such classes,
\begin{equation}\label{eq:hardclass}
\mathcal{H} \;=\; \{1,\ 121,\ 169,\ 289,\ 361,\ 529\} \pmod{840},
\end{equation}
being $1^2, 11^2, 13^2, 17^2, 19^2, 23^2$. All six lie in $1\pmod{24}$, but that coarse
class contains $24$ reduced classes modulo $840$, only six of which are squares. The
other three quarters have already been settled, yet their early failed rungs would
still enter the denominator of any reported failure share.

\begin{observation}[the easy primes contribute only level J]\label{obs:dilution}
Through $10^{11}$, the censuses of $(1\bmod24)$ and $\mathcal{H}$ have identical depth
counts at every rung $r\ge11$ and identical level-RC, level-S, and level-R totals. Every
extra failure from $(1\bmod24)\setminus\mathcal{H}$ occurs at $r=3$ and has a Jacobi
certificate. Restricting to $\mathcal{H}$ therefore removes only already-explained
failures and preserves the entire residual and deep-ladder population in the overlap.
\end{observation}

Including the coarse classes thus understates the residual share. All subsequent depth
and failure-share statistics use $\mathcal{H}$.

\begin{observation}[the levels on $\mathcal{H}$]\label{obs:hard13}
The classified census on the hard class is:
\begin{center}
\small
\begin{tabular}{lrrrrrr}
\toprule
& primes & level J & level RC & level S & level R & residual share \\
\midrule
$p < 10^6$    & $2{,}370$          & $3{,}318$         & $2$      & $0$     & $112$          & $3.2634\%$ \\
$p < 10^7$    & $20{,}513$         & $24{,}376$        & $2$      & $0$     & $662$          & $2.6438\%$ \\
$p < 10^8$    & $179{,}468$        & $186{,}590$       & $4$      & $2$     & $4{,}018$      & $2.1079\%$ \\
$p < 10^9$    & $1{,}587{,}581$    & $1{,}473{,}032$   & $18$     & $4$     & $26{,}317$     & $1.7552\%$ \\
$p < 10^{10}$ & $14{,}215{,}707$   & $11{,}942{,}520$  & $87$     & $8$     & $181{,}523$    & $1.4972\%$ \\
$p < 10^{11}$ & $128{,}671{,}219$  & $99{,}140{,}159$  & $454$    & $47$    & $1{,}302{,}101$& $1.2964\%$ \\
$p < 10^{12}$ & $1{,}175{,}215{,}396$ & $838{,}863{,}859$ & $2{,}279$ & $242$ & $9{,}664{,}186$ & $1.1389\%$ \\
$p < 10^{13}$ & $10{,}814{,}395{,}686$ & $7{,}211{,}448{,}416$ & $12{,}112$ & $1{,}309$ & $73{,}736{,}922$ & $1.0121\%$ \\
$p < 10^{14}$ & $100{,}154{,}081{,}439$ & $62{,}821{,}444{,}770$ & $66{,}767$ & $7{,}373$ & $576{,}087{,}976$ & $0.9087\%$ \\
$p < 10^{15}$ & $932{,}642{,}160{,}749$ & $553{,}373{,}502{,}288$ & $386{,}596$ & $43{,}257$ & $4{,}590{,}145{,}833$ & $0.8227\%$ \\
\bottomrule
\end{tabular}
\end{center}
Two features stand out.
\begin{itemize}
\item \textbf{Level S occurs only at $r = 51$}, across all $43{,}257$ instances, while
its sibling level RC spreads to
$r = 51$ ($382{,}218$), $75$ ($4{,}137$), $91$ ($169$), $99$ ($59$), $115$ ($7$) and
$123$ ($6$). The rung grid does not force this localization --- but
Corollary~\ref{cor:Srungs} does: $51$ is the smallest rung at which level S is possible
at all, and $75$, $91$, $99$, $115$ are rungs at which it is impossible. Each of the
$43{,}257$ therefore has support subgroup $H = \langle 2 \bmod 51 \rangle$, the unique
admissible one. The appearance of level RC at $r = 123$ sharpens the reachability
point: $123$ is on the admissible list of Corollary~\ref{cor:Srungs} and is now
demonstrably reached carrying local obstructions, yet level S still declines to appear
there --- at $r=123$ the admissible subgroup must arise as an actual support subgroup,
which evidently remains rare.
\item \textbf{The residual share falls monotonically}, from $3.26\%$ to $0.8227\%$,
while the residual count rises from $112$ to $4{,}590{,}145{,}833$. The proportions and
the counts therefore tell different stories.
\end{itemize}
\end{observation}

\subsection{The central empirical thesis}\label{sec:thesis}

We state the thesis the census supports, and---as carefully---the stronger statements it
does not.

\begin{quote}
\emph{Local certificates asymptotically dominate the observed failure landscape; the
remaining risk is concentrated in a sparse, divisor-poor tail.}
\end{quote}

Formally we separate three hypotheses, which are often conflated and have very
different status.

\begin{description}
\item[H1 (local dominance).] Among failed rungs \emph{of primes in $\mathcal{H}$}, the
residual share tends to zero:
\[
\Pr\bigl(\text{level R} \,\big|\, \text{failure},\ p \le X\bigr) \longrightarrow 0
\qquad (X \to \infty).
\]
This is the hypothesis the data speaks to directly. \emph{It does not imply the
Erd\H{o}s--Straus conjecture}: a vanishing share is compatible with a growing number of
residual misses and says nothing about whether one prime can miss every rung.
\item[H2 (divisor threshold).] For each rung $r$ there is a threshold $f(r)$ with
$\tau(q^2) \ge f(r) \Rightarrow N_r(p) > 0$ whenever the target class is reachable.
Residual misses still occur at $\tau(q^2)=81$, so no small constant threshold is
visible. A threshold depending on $r$ remains open, and Section~\ref{sec:occupancy}
shows why naive occupancy cannot establish one.
\item[H3 (unbounded but slow depth).] There is no uniform constant $R$ with
$D(p)\le R$ for every $p$, but the record envelope
\[
M(X)=\max_{p<X,\ p\in\mathcal H}D(p)
\]
grows exceptionally slowly. The move from $107$ through $131$ to $155$ after long
plateaus is
consistent with this hypothesis and inconsistent with treating the old record as a
natural barrier. Finite computation cannot distinguish genuine unboundedness from a
larger constant bound, so H3 is deliberately stated as the numerical expectation, not
as a conclusion.
\end{description}

\subsection{Evidence for H1}\label{sec:h1}

The cumulative residual share on $\mathcal{H}$ is
\begin{equation}\label{eq:hseries}
3.26\%,\ 2.64\%,\ 2.11\%,\ 1.76\%,\ 1.50\%,\ 1.30\%,\ 1.14\%,\ 1.01\%,\ 0.9087\%,\ 0.8227\%
\qquad (X = 10^6,\ldots,10^{15}).
\end{equation}
This decline is monotone, but cumulative censuses are nested and therefore unsuitable
for a clean holdout test. Differencing them into disjoint decades gives
\[
\begin{array}{lccccccccc}
[10^{k-1}, 10^k) & k=7 & 8 & 9 & 10 & 11 & 12 & 13 & 14 & 15 \\
\text{residual share} & 2.5454\% & 2.0269\% & 1.7038\% & 1.4608\% & 1.2688\% &
1.1178\% & 0.9954\% & 0.8953\% & 0.8116\%
\end{array}
\]
which is also monotone. A log-power fit to the first six decades predicts $0.9838\%$
for $[10^{12},10^{13})$, against the observed $0.9954\%$, a $1.2\%$ relative error on
a disjoint range; the same family fitted through $10^{13}$ predicts $0.8844\%$ and
$0.7973\%$ for the two decades that were computed only afterwards, against the observed
$0.8953\%$ and $0.8116\%$ ($1.2\%$ and $1.8\%$ relative error). This supports H1,
but it does not establish a law or justify remote extrapolation. Most importantly, the
falling share coexists with a rapidly increasing residual count.

\subsection{Against random-divisor models: an occupancy test}\label{sec:occupancy}

There is a tempting heuristic for H2: if the $\tau(q^2)$ divisors were spread evenly over
the $|H|$ reachable classes, the chance of missing one target class would be about
$e^{-m}$ with $m = \tau(q^2)/|H|$, so a large lattice would force a hit. We tested this
directly on the recorded $p\equiv1\pmod{24}$ rung sample below $2 \times 10^8$, retaining
only rungs whose target class is reachable
(levels HIT and R), binning by $m$ and comparing the observed miss rate to $e^{-m}$
evaluated at the \emph{lower} edge of each bin---the most generous value the null takes
anywhere in the bin.

\begin{table}[ht]\centering
\caption{Observed residual-miss rate against the random-occupancy null on the recorded
$p\equiv1\pmod{24}$ sample, $p < 2\times10^8$.}
\label{tab:occ}
\small
\begin{tabular}{lrrrr}
\toprule
$m = \tau(q^2)/|H|$ & rungs & observed miss rate & $\max e^{-m}$ in bin & excess \\
\midrule
$[1,2)$   & $3{,}304$   & $0.4837$   & $3.7 \times 10^{-1}$  & $1.3\times$ \\
$[2,4)$   & $5{,}543$   & $0.4415$   & $1.4 \times 10^{-1}$  & $3.3\times$ \\
$[4,8)$   & $44{,}730$  & $0.0194$   & $1.8 \times 10^{-2}$  & $1.1\times$ \\
$[8,16)$  & $397{,}224$ & $3.8 \times 10^{-3}$ & $3.4 \times 10^{-4}$ & $11\times$ \\
$[16,32)$ & $130{,}304$ & $1.0 \times 10^{-3}$ & $1.1 \times 10^{-7}$ & $9.2 \times 10^{3}\times$ \\
$[32,64)$ & $425{,}599$ & $9.4 \times 10^{-6}$ & $1.3 \times 10^{-14}$ & $7.1 \times 10^{8}\times$ \\
\bottomrule
\end{tabular}
\end{table}

The null is not merely inaccurate but wrong by growing orders of magnitude exactly where
H2 would need it: at $m \in [16,32)$ residual misses are nearly ten thousand times
commoner than independence predicts, and beyond that the null predicts essentially none
while misses still occur. Normalizing by $\varphi(r)$ instead of $|H|$ changes the table
hardly at all, so this is \emph{not} subgroup confinement---the support subgroup is
usually everything. It is genuine dependence among the divisors of $q^2$ modulo $r$,
which is unsurprising on reflection: those divisors are $\prod \ell_i^{f_i}$ with
$f_i \le 2e_i$, a rigid multiplicative lattice, not a random set.

The consequence for \S\ref{sec:doorways} is concrete: any route to H2 through
equidistribution of divisors in residue classes---on average, or for typical
integers---will not do, because the relevant configurations are precisely the atypical
ones. This sharpens Direction III from ``hard'' to ``hard for an identifiable reason.''

\subsection{How not to model the ladder}\label{sec:model}

The rungs of one prime are not independent trials. They share $p$, while the shifted
values $(p+3)/4,(p+7)/4,\ldots$ are consecutive integers with correlated
factorizations. A statistical model must therefore treat the ladder as a
rung-dependent survival process, not as repeated Bernoulli sampling. The observed
shares are descriptive marginals; confidence intervals based on independence would be
misleading.

Residual misses are nevertheless consistently divisor-poor: their median
$\tau(q^2)$ is $27$, compared with $81$ at hitting rungs in the matched summaries. This
is a correlation, not a consequence of \eqref{eq:fourier}; a small main term makes a
zero possible, not probable.

The resulting decomposition of the covering problem is:

\begin{itemize}
\item a \emph{local layer}, entirely finite: which rungs carry a real-character or
support certificate is a computation in $(\mathbb{Z}/r)^\times$ given the residues of
the primes of $q$;
\item a \emph{residual layer}: prove that level-R misses cannot occupy every rung of
some prime. This is the content the conjecture actually needs, and H1 does not supply
it---a vanishing share of failures is not a statement about any individual ladder.
\S\ref{sec:occupancy} shows the natural probabilistic route is blocked, and
\S\ref{sec:doorways} that the natural analytic ones are too.
\end{itemize}

\section{The computation}\label{sec:computation}

The census requires the complete factorization of every shifted value
$A=(p+r)/4$ encountered before the first hit. Factoring these values independently is
unnecessarily expensive; their arithmetic-progression structure permits a bulk sieve.

\subsection{Bulk sieve-factorization}\label{sec:bulk}

Within one of the six classes modulo $840$, write $p=p_0+840j$. At a fixed rung,
\begin{equation}\label{eq:apA}
A=\frac{p+r}{4}=A_0+210j, \qquad A_0=\frac{p_0+r}{4}.
\end{equation}
For every prime $\ell\nmid210$,
\[
\ell\mid A \quad\Longleftrightarrow\quad
j\equiv-A_0\,\overline{210}\pmod\ell.
\]
Thus each prime $\ell$ visits exactly the entries it divides. Sieving through
$\sqrt{A_{\max}}$ leaves a prime cofactor, eliminating primality tests and Pollard rho
from the main path. Once only a thin set of primes remains on the ladder, the scanner
switches back to per-value factorization. Against the naive implementation this gives
speedups of $7.3\times$, $12.8\times$, and $20.3\times$ at $10^{10}$, $10^{11}$, and
$10^{12}$, respectively, with identical output.

\subsection{Scaling}\label{sec:cost}

The measured per-prime cost grows slowly with height, so total cost is dominated by the
number of primes. The range splits into independent shards and merges by adding
histograms. On the machine used for the final run, the complete $10^{13}$ hard-class
census took $17$ minutes on $12$ physical cores. The extension to $10^{15}$ was
produced by a CUDA implementation of the same pipeline on a single NVIDIA L4 GPU in
$2.6$ hours ($\approx 9$\,ns per prime); its statistical output and its certificate
set are byte-identical to the CPU reference on every range where both were run, which
is the standard its adoption required. Different window sizes, shard decompositions,
and the CPU/GPU implementation pair all produce identical statistics.

\subsection{Verification protocol}\label{sec:protocol}

The optimized scanners (CPU and GPU) are checked against an independent
exact-arithmetic reference and against the earlier per-value implementation. Repeated
runs, thread counts, window sizes, and independent shardings give identical counts.
The final shards ($192$ at $10^{13}$, $400$ at $10^{15}$) are checked for gaps and
overlaps, and every fine-certificate and deep-prime record---$20{,}001$ at $10^{13}$,
$585{,}677$ at $10^{15}$---is re-derived in exact arithmetic on a machine different
from the one that produced it. The verifier additionally enforces
Proposition~\ref{prop:Spossible} mechanically: a level-S claim at an inadmissible
rung, or one whose support subgroup omits $4$, is rejected as falsifying rather than
merely erroneous. Each reported solution is verified by the integer identity
$4ABC=p(AB+BC+CA)$.

\section{Relation to prior work}\label{sec:prior}

The divisor-congruence criterion of Lemma~\ref{lem:divisor} is classical and appears in
several equivalent guises. Mordell's identities \cite{Mordell1969} and the modular
equations catalogued by Salez \cite{Salez2014} are, in the present language, individual
rungs: a polynomial identity valid on a residue class of $p$ is a rule for choosing $A$
(hence $r$) together with a divisor $u$ that provably lands in the target class.
Yamamoto \cite{Yamamoto1965} and Schinzel \cite{Schinzel2000} established the
quadratic obstruction that Corollary~\ref{cor:blocking} localizes to a single rung.
Elsholtz and Tao \cite{ElsholtzTao2013} give the Type I / Type II taxonomy used in
Proposition~\ref{prop:types}, together with mean-value theorems for the number of
solutions; their framework subsumes the parametrization here. Vaughan \cite{Vaughan1970}
bounded the exceptional set by $\ll N\exp(-c(\log N)^{2/3})$. The conjecture is
verified numerically to $10^{18}$ \cite{MihneaBogdan2025}, after Swett \cite{Swett1999}
and Salez \cite{Salez2014}; a recent preprint \cite{Unified2026} obtains density one
within $n \equiv 1 \pmod 4$ using divisor conditions rather than congruence identities,
and is the closest antecedent to the perspective taken here.

Accordingly we do \emph{not} claim the divisor-lattice viewpoint as a new framework. Its
ingredients are known; what we add is the certificate hierarchy of
Definition~\ref{def:hierarchy}, its exhaustive measurement, and the released
certificates.

\section{Directions: Fourier interference and its obstacles}\label{sec:doorways}

The Fourier analogy becomes exact after centering the divisor lattice. Write
$q=\prod_\ell \ell^{e_\ell}$ and let $G_r=(\mathbb Z/r)^\times$. A divisor
$u=\prod_\ell\ell^{j_\ell}$ of $q^2$ determines, in $G_r$, the centered point
$u/q=\prod_\ell\ell^{k_\ell}$ with $k_\ell=j_\ell-e_\ell$ and
$-e_\ell\leq k_\ell\leq e_\ell$. Thus the moving target $-q$ becomes the fixed target
$-1$. Define the multiplicity measure
\[
 \nu_{q,r}(g)=\#\left\{(k_\ell): -e_\ell\leq k_\ell\leq e_\ell,
       \ \prod_{\ell\mid q}\ell^{k_\ell}=g\text{ in }G_r\right\}.
\]
Then $N_r(p)=\nu_{q,r}(-1)$. With the convention
$\widehat\nu(\chi)=\sum_{g\in G_r}\nu(g)\chi(g)$, and writing
$\chi(\ell)=e^{i\theta_{\chi,\ell}}$, its finite Fourier transform factors as
\begin{equation}\label{eq:centered-transform}
 \widehat\nu_{q,r}(\chi)
 =\prod_{\ell^{e_\ell}\parallel q}D_{e_\ell}(\theta_{\chi,\ell}),
 \qquad
 D_e(\theta)=\sum_{k=-e}^{e}e^{ik\theta}
 =\frac{\sin((2e+1)\theta/2)}{\sin(\theta/2)},
\end{equation}
where the quotient has the limiting value $2e+1$ at $\theta=0$. Fourier inversion
therefore rewrites Theorem~\ref{thm:fourier} as
\begin{equation}\label{eq:centered-inversion}
 N_r(p)=\frac{1}{\varphi(r)}\left\{
 \tau(q^2)+\sum_{\chi\neq\chi_0}\chi(-1)
 \prod_{\ell^{e_\ell}\parallel q}D_{e_\ell}(\theta_{\chi,\ell})\right\}.
\end{equation}
Every local factor is real, and $\chi(-1)=\pm1$. A miss is consequently literal
destructive interference: the nonprincipal spectrum cancels the positive mass
$\tau(q^2)$ exactly.

This gives a deterministic, if not yet uniform, non-cancellation certificate. Put
\begin{equation}\label{eq:interference-mass}
 \mathfrak I_r(p)=\sum_{\chi\neq\chi_0}
 \prod_{\ell^{e_\ell}\parallel q}
 \frac{|D_{e_\ell}(\theta_{\chi,\ell})|}{2e_\ell+1}.
\end{equation}
\begin{proposition}[spectral non-cancellation]\label{prop:spectral-noncancellation}
The triangle inequality in \eqref{eq:centered-inversion} gives
\begin{equation}\label{eq:spectral-lower}
 N_r(p)\ \geq\ \frac{\tau(q^2)}{\varphi(r)}\bigl(1-\mathfrak I_r(p)\bigr),
 \qquad
 \mathfrak I_r(p)<1\ \Longrightarrow\ N_r(p)>0.
\end{equation}
\end{proposition}
\begin{proof}
Each normalized product in \eqref{eq:interference-mass} is the absolute value of the
corresponding nonprincipal term in \eqref{eq:centered-inversion}, divided by
$\tau(q^2)$. Subtract their sum from the principal term.
\end{proof}
Unlike a random-occupancy model, \eqref{eq:spectral-lower} is exact and retains the
dependence among divisors. Its weakness is also precise: it discards the signs between
nonprincipal modes.

The hierarchy now has a spectral reading. Let
\[
 B_{q,r}=\prod_{\ell^{e_\ell}\parallel q}
 \{\ell^{-e_\ell},\ldots,1,\ldots,\ell^{e_\ell}\}\subseteq G_r .
\]
Then $N_r(p)>0$ exactly when $-1\in B_{q,r}$. A J- or RC-certificate is a maximally
coherent odd real mode, trivial on every generator of the support subgroup $H$. At
level S the same obstruction is detected by an odd higher-order character in the
annihilator of $H$; it is invisible only if one restricts attention to real characters.
At level R, $-1\in H$ but the finite exponent box still misses it, so no annihilator
character separates the target: cancellation is spread among partially damped modes.
Equivalently, this is a product-set problem. Finite-group product-set estimates turn
sufficiently small expansion of the successive factors in $B_{q,r}$ into a stabilizer
subgroup, while expansion without a stabilizer is the combinatorial counterpart of
Fourier damping. This structure-versus-mixing dichotomy is the main line of evidence
behind the interference picture.

\paragraph{Direction I: spectral mixing on one rung.}
The immediate aim is to control \eqref{eq:interference-mass} from the orders of the
residues $\ell\bmod r$ and the exponents $e_\ell$. If a nonprincipal character is
nontrivial on several prime factors of $q$, the corresponding normalized Dirichlet
kernels are damped. One needs enough such damping collectively to make
$\mathfrak I_r(p)<1$, or a signed refinement of \eqref{eq:spectral-lower} when the
absolute-value bound is too expensive. Product-set growth offers a dual route: prove
that $B_{q,r}$ either has a proper stabilizer, already recorded by the hierarchy, or
covers $-1$ once its exponent box is sufficiently rich. This formulation explains why
$\tau(q^2)$ alone cannot be a threshold: the phases and their multiplicities matter.

\paragraph{Direction II: decorrelation across the ladder.}
For $q_r=p(p+r)/4$, let $E_r(p)$ denote the nonprincipal part of
\eqref{eq:centered-inversion}, including the factor $1/\varphi(r)$, and set
$\mathcal R_p(R)=\{r\leq R:r\equiv3\pmod4,\ \gcd(q_r,r)=1\}$. Then
\begin{equation}\label{eq:ladder-interference}
 \sum_{r\in\mathcal R_p(R)}N_r(p)
 =\sum_{r\in\mathcal R_p(R)}
   \frac{\tau(q_r^2)}{\varphi(r)}
 +\sum_{r\in\mathcal R_p(R)}E_r(p).
\end{equation}
(A noncoprime rung already hits by Lemma~\ref{lem:coprime}.) Hence total interference
cannot persist through rung $R$ if
\[
 \left|\sum_{r\in\mathcal R_p(R)}E_r(p)\right|
 <\sum_{r\in\mathcal R_p(R)}\frac{\tau(q_r^2)}{\varphi(r)}.
\]
This is the rigorous form of the proposed ``no total interference'' mechanism. The
difficulty is that both the group and the signal change with $r$: the signal at rung
$r$ is built from the factorization of $q_r$, so a standard large-sieve estimate does
not apply unchanged. The sought-for statement is a cross-rung uncertainty principle
for these arithmetically coupled spectra.

\paragraph{Direction III: arithmetic input for the moving signal.}
All analytic routes meet the same shifted-factorization obstacle. The Frobenius class
of $p$ controls $\chi(p)$, but not the prime factors of $(p+r)/4$. Selberg--Delange
theory gives useful models for the frequency with which all those factors lie in a
fixed subgroup---for an index-two condition, a count of order
$c_rx(\log x)^{-1/2}$ follows from the usual branch-point argument
\cite[Ch.~II.5]{Tenenbaum2015}---but conditioning on $p$ prime turns the question into a
multiplicative correlation along a linear form. Chebotarev alone does not see that
correlation, and sieve methods encounter the parity barrier.

For residual modes, the natural input is divisor distribution: Hooley's
$\Delta$-function \cite{Hooley1979}, Hall--Tenenbaum \cite{HallTenenbaum1988}, and
Ford's work on divisors in prescribed ranges \cite{Ford2008}. Those theories control
typical integers; here one needs uniformity for the structured sequence
$q_r^2=p^2((p+r)/4)^2$, ultimately for every $p$. No known method makes that jump. The
Fourier formulation does not remove this obstacle, but identifies the quantities that
must be controlled: local phase damping on one rung or signed spectral decorrelation
across several rungs.

\section{Problems}\label{sec:problems}

\begin{problem}[spectral census]\label{prob:sd}
Augment the census with the normalized local factors in
\eqref{eq:interference-mass}, the signed error $E_r(p)$, and the margin
$1-\mathfrak I_r(p)$. Determine how often \eqref{eq:spectral-lower} certifies a hit,
which characters dominate residual misses, and whether every computed ladder contains
a rung with $\mathfrak I_r(p)<1$. For coherent J, RC and S modes, compare the observed
frequencies with Selberg--Delange models that retain the coupling between $p$ and
$(p+r)/4$.
\end{problem}

\begin{problem}[structure versus mixing]\label{prob:divisors}
Prove a nontrivial condition on the orders of $\ell\bmod r$ and the exponents $e_\ell$
under which $-1\in H$ forces $\mathfrak I_r(p)<1$. Equivalently, give a product-set
theorem forcing $-1\in B_{q,r}$ once stabilizer obstructions have been excluded.
Even a result for prime $r$, or for squarefree $q$, would separate genuine subgroup
confinement from finite-box interference. A bound depending only on $\tau(q^2)$ cannot
capture this distinction.
\end{problem}

\begin{problem}[cross-rung non-cancellation]\label{prob:interference}
Find a range $R=R(p)<p$ on which the signed spectral errors in
\eqref{eq:ladder-interference} satisfy
\[
 \left|\sum_{r\in\mathcal R_p(R)}E_r(p)\right|
 <\sum_{r\in\mathcal R_p(R)}\frac{\tau(q_r^2)}{\varphi(r)}.
\]
Any such estimate for every $p\in\mathcal H$ proves that some rung hits. An average or
almost-all version would not prove the conjecture, but would test whether the observed
slow growth of record depth is explained by decorrelation rather than by a fixed bound.
Identify precisely what replacement for the large sieve can accommodate the moving
factorizations $q_r=p(p+r)/4$.
\end{problem}

\begin{problem}\label{prob:bound}
Determine whether $M(X)=\max_{p<X,\ p\in\mathcal H}D(p)$ is unbounded. Short of that,
derive a tail estimate for $D(p)$ strong enough to distinguish a slowly diverging record
from a fixed ceiling, and compare record setters with the smallest spectral margin
$\min_{r<D(p)}|1-\mathfrak I_r(p)|$ on their failed rungs. A bound $D(p)=o(p)$ for
every sufficiently large $p$ would already prove the conjecture on $\mathcal H$ after
finite verification, so it is not a modest intermediate target.
\end{problem}

\begin{problem}\label{prob:tail}
Corollary~\ref{cor:Srungs} says which rungs \emph{can} carry a level-S failure; it does
not say which do. Below $10^{15}$ only $r = 51$ is occupied, though $r = 119$ and $123$ are
admissible and are visited by every ladder deeper than them --- such ladders exist (the
record is $155$), and level RC now demonstrably reaches $123$. Estimate the density of primes whose support subgroup at rung $119$ or
$123$ equals the unique admissible subgroup there, and predict the height at which a second
occupied rung should first appear. A prediction that survives one further decade would turn
Corollary~\ref{cor:Srungs} from a constraint into a model.
\end{problem}

\section{Conclusion}\label{sec:conclusion}

The certificate hierarchy explains almost every failed rung in the hard class, usually
through the Jacobi symbol alone. What remains is small in proportion but not in count:
the residual share falls to $0.8227\%$, while $4.59$ billion residual misses remain
below $10^{15}$. Those misses occur disproportionately at divisor-poor shifted values,
yet their divisors are far too dependent modulo $r$ for a random-occupancy argument. The
one level that no real-character certificate detects is not scattered either: because
$4A = p + r$ puts $4$ in every support subgroup, level S can occur at only ten rungs
below $300$, and every observed instance sits at the smallest of them.

The depth data carries a separate message. The maximum first-hit rung stayed at $107$
for a long interval and then rose through $131$ to $155$, with intermediate depths
populated. This is
the pattern expected from a very thin extreme tail. It provides no evidence for a
uniform $O(1)$ depth bound. Our numerical expectation is instead that the record depth
is unbounded but grows extraordinarily slowly; the census is far too short to identify
that growth law.

None of this proves the conjecture. A vanishing share of residual failures does not
prevent one prime from missing every rung, and an observed record is neither an upper
bound nor a theorem of divergence. The useful outcome is the separation itself: local
obstructions are finite and certifiable, while the remaining problem concerns the
structured divisors of $p^2((p+r)/4)^2$ along shifted prime values. That is the part a
proof must control.

\section*{Data and code availability}\label{sec:data-availability}
\addcontentsline{toc}{section}{Data and code availability}

The complete public repository is
\url{https://github.com/Attitude-Adjuster/erdos-strauss}. Questions, corrections, code
problems, and independent reproduction reports may be submitted at
\url{https://github.com/Attitude-Adjuster/erdos-strauss/issues}.

The repository contains the exact-arithmetic reference implementation, the
independent C\raisebox{0.2ex}{\small ++} scanners (the frozen reference and an
optimized variant proved output-identical against it), the CUDA scanner, the
certificate verifier---including the mechanical check of
Proposition~\ref{prop:Spossible}---the per-shard summaries of both censuses ($192$
shards at $10^{13}$, $400$ at $10^{15}$), the merged depth and level histograms, the
$20{,}001$ fine-certificate and deep-prime records of the $10^{13}$ census, the
$585{,}677$ of the $10^{15}$ census, and the verification transcripts. The analysis
scripts reproduce the cumulative and disjoint-decade statistics. Level-J failures are
summarized rather than stored individually because each is regenerated directly from
$(p,A,r)$.

Each solution certificate is a triple $(p,A,u)$ and is checked through
$4ABC=p(AB+BC+CA)$. The per-shard summaries are published so that small pieces of the
census can be reproduced independently rather than requiring the full run.

\begin{table}[ht]\centering
\caption{Numerical check of \eqref{eq:fourier} against direct enumeration, reporting
both the correct (conjugated) and the mis-transcribed (unconjugated) evaluations; the
latter counts the class $(-q)^{-1}$ instead of $-q$.}\label{tab:verify}
\begin{tabular}{rrrlrrr}
\toprule
$p$ & $A$ & $r$ & rung status & \eqref{eq:fourier} & unconjugated & direct \\
\midrule
   73 &   19 &  3 & fails (J)   & $0$ & $0$ & 0 \\
   73 &   20 &  7 & hits        & $8$ & $6$ & 8 \\
   73 &   21 & 11 & hits        & $4$ & $2$ & 4 \\
   73 &   23 & 19 & fails (J)   & $0$ & $0$ & 0 \\
 1129 &  285 & 11 & hits        & $8$ & $8$ & 8 \\
21169 & 5293 &  3 & fails (J)   & $0$ & $0$ & 0 \\
21169 & 5300 & 31 & hits        & $6$ & $6$ & 6 \\
\bottomrule
\end{tabular}
\end{table}

\subsection*{Acknowledgement}
A detailed critical reading identified the composite-$r$ gap in the original
obstruction statement and sharpened the discussion of the analytic obstacles. It also
supplied the independent divisor counts used in Table~\ref{tab:verify}.

\begin{thebibliography}{99}

\bibitem{Davenport2000} H.~Davenport,
\emph{Multiplicative Number Theory}, 3rd ed., rev.\ H.~L.~Montgomery,
Graduate Texts in Mathematics \textbf{74}, Springer, 2000.

\bibitem{Erdos1950} P.~Erd\H{o}s,
\emph{Az $\frac{1}{x_1} + \cdots + \frac{1}{x_n} = \frac{a}{b}$ egyenlet eg\'esz
sz\'am\'u megold\'asair\'ol}, Mat.\ Lapok \textbf{1} (1950), 192--210.

\bibitem{ElsholtzTao2013} C.~Elsholtz and T.~Tao,
\emph{Counting the number of solutions to the Erd\H{o}s--Straus equation on unit
fractions}, J.~Aust.\ Math.\ Soc.\ \textbf{94} (2013), 50--105.

\bibitem{Ford2008} K.~Ford,
\emph{The distribution of integers with a divisor in a given interval},
Ann.\ of Math.\ \textbf{168} (2008), 367--433.

\bibitem{HallTenenbaum1988} R.~R.~Hall and G.~Tenenbaum,
\emph{Divisors}, Cambridge Tracts in Mathematics \textbf{90}, CUP, 1988.

\bibitem{Hooley1979} C.~Hooley,
\emph{On a new technique and its applications to the theory of numbers},
Proc.\ London Math.\ Soc.\ \textbf{38} (1979), 115--151.

\bibitem{MihneaBogdan2025} S.~Mihnea and D.~C.~Bogdan,
\emph{Further verification and empirical evidence for the Erd\H{o}s--Straus
conjecture}, arXiv:2509.00128 (2025).

\bibitem{Mordell1969} L.~J.~Mordell,
\emph{Diophantine Equations}, Academic Press, 1969.

\bibitem{Salez2014} S.~E.~Salez,
\emph{The Erd\H{o}s--Straus conjecture: new modular equations and checking below
$10^{17}$}, arXiv:1406.6307 (2014).

\bibitem{Schinzel2000} A.~Schinzel,
\emph{On sums of three unit fractions with polynomial denominators},
Funct.\ Approx.\ Comment.\ Math.\ \textbf{28} (2000), 187--194.

\bibitem{Swett1999} A.~Swett,
\emph{The Erd\H{o}s--Straus conjecture}, electronic verification to $10^{14}$, 1999--2011.

\bibitem{Tenenbaum2015} G.~Tenenbaum,
\emph{Introduction to Analytic and Probabilistic Number Theory}, 3rd ed.,
Graduate Studies in Mathematics \textbf{163}, AMS, 2015.

\bibitem{Unified2026}
\emph{A unified parametric approach to the Erd\H{o}s--Straus conjecture with explicit
solutions for a set of integers of natural density one}, arXiv:2602.20036 (2026).

\bibitem{Vaughan1970} R.~C.~Vaughan,
\emph{On a problem of Erd\H{o}s, Straus and Schinzel},
Mathematika \textbf{17} (1970), 193--198.

\bibitem{Yamamoto1965} K.~Yamamoto,
\emph{On the Diophantine equation $\frac{4}{n} = \frac{1}{x}+\frac{1}{y}+\frac{1}{z}$},
Mem.\ Fac.\ Sci.\ Kyushu Univ.\ Ser.~A \textbf{19} (1965), 37--47.

\end{thebibliography}

\end{document}
